% =========================================================
% SECCIÓN 3.8 - CONCLUSIONES DEL CAPÍTULO
% =========================================================


Este capítulo ha desarrollado la generalización del modelo Dynamic Energy Budget (DEB) de \parencite{eriksenDynamicModellingFeed2022} desde la escala individual larval hasta la escala de reactor batch de \SI{20}{\litre}, formulando un sistema de cinco ecuaciones diferenciales ordinarias que gobiernan la dinámica del sustrato orgánico, la reserva de asimilados, la biomasa estructural, los lípidos de reserva y la acumulación de \ch{CO2}. Los principales hallazgos y decisiones de modelado se sintetizan a continuación.

\paragraph{Arquitectura de compartimentos y flujos.} Se adoptó un sistema de cinco compartimentos de carbono interconectados —sustrato $S$, reserva $E$, biomasa estructural $B$, lípidos $L$ y \ch{CO2} acumulado $C_{\ch{CO2}}$— que permite rastrear la destinación del carbono ingerido desde la ingesta hasta la disipación gaseosa. Los flujos metabólicos —ingesta, asimilación, mantenimiento, síntesis de biomasa, síntesis de lípidos y disipación— se derivaron directamente del marco DEB estándar, preservando la interpretabilidad bioenergética de cada parámetro.

\paragraph{Formulación del sistema de EDO.} Las ecuaciones de balance resultantes constituyen un sistema acoplado no lineal cuya estructura garantiza la conservación de masa de carbono en régimen aeróbico. La hipótesis de estado cuasi-estacionario para la reserva de asimilados ($\mathrm{d}E/\mathrm{d}t \approx 0$) se justificó por separación de escalas temporales entre el amortiguamiento metabólico (horas) y el crecimiento somático (días). La producción de \ch{CO2} se expresó como suma de los costos de mantenimiento y los \emph{overheads} de síntesis, vinculándose con el flujo másico medido por los sensores NDIR del Capítulo \ref{chap:instrumentacion} mediante la conversión molar $M_{\ch{CO2}}/M_{\ch{C}}$.

\paragraph{Cinética de asimilación y modulación ambiental.} La ingesta larval se modeló mediante una funcional respuesta tipo Holling II, que captura la saturación por disponibilidad de sustrato y escala isométricamente con la superficie corporal ($B^{2/3}$). La dependencia térmica se incorporó mediante el modelo de Sharpe--Schoolfield con inactivación de alta temperatura, parametrizado con semillas entomológicas ($T_A \in [2000, 8000]$ K, $T_{\mathrm{opt}} \approx 303$ K) y dejado como variable libre de calibración local. La humedad relativa no se incluyó como variable cinética explícita, dado que el protocolo experimental la mantiene dentro del rango óptimo ($60\%$--$80\%$); las desviaciones transitorias se absorberán en la calibración de $\tau_h$ y $a_{\max}$ en el Capítulo \ref{chap:hibrido}.

\paragraph{Parámetros y restricciones físicas.} Se distinguieron parámetros fijos adoptados de la literatura ($\kappa_A, a_{\max}, m, Y_B, Y_L, \kappa$) de variables libres de calibración local ($\tau_h, \beta_0, \rho_{\mathrm{conv}}, T_A, T_{\mathrm{opt}}, T_H$), con rangos semilla y restricciones de caja que garantizan eficiencias termodinámicas menores al $100\%$ y conservación de masa. El factor de conversión $\rho_{\mathrm{conv}}$ vincula la biomasa húmeda destructiva con el carbono estructural, puente esencial para la confrontación modelo--sensor.

Se identificó además el \emph{subconjunto de parámetros de primera sensibilidad} respecto a la pendiente de emisión de \ch{CO2} en ventanas cerradas: $\{\tau_h, \rho_{\mathrm{conv}}, T_{\mathrm{opt}}\}$. Estos tres parámetros modulan directamente la ingesta, el mantenimiento y la tasa metabólica en las condiciones del experimento, mientras que $\beta_0$, $T_A$ y $T_H$ operan en régimen de segundo orden en el rango térmico controlado. Esta jerarquía de sensibilidad será explotada en el Capítulo \ref{chap:hibrido} para reducir el espacio de corrección del calibrador neuronal y preservar la interpretabilidad biofísica.

\paragraph{Análisis estructural.} Se verificó la consistencia dimensional de todos los flujos (g C/min), el cierre del balance de carbono en régimen aeróbico y el comportamiento cualitativo biológicamente plausible: no negatividad del dominio $\Omega$, monotonicidad decreciente del sustrato y acotación superior de la biomasa por el carbono total disponible. Las demostraciones formales de existencia, unicidad y estabilidad de las soluciones, así como el análisis de sensibilidad paramétrica, se remitieron al Anexo \ref{anexo:formulacion_matematica} para mantener la fluidez del capítulo.

\paragraph{Del modelo mecanicista al sistema híbrido.} El modelo DEB formulado en este capítulo constituye la línea base bioenergética del sistema híbrido: predice cuánto \ch{CO2} debería emitirse bajo condiciones de aerobiosis estricta y homogeneidad del sustrato. Sin embargo, la biomasa estructural $B(t)$ —variable de estado crítica para el cálculo del mantenimiento y la ingesta— no es medible en tiempo real sin perturbación destructiva del sistema. Esta imposibilidad instrumental constituye la brecha que el Capítulo \ref{chap:hibrido} resuelve mediante un estimador híbrido modelo--datos.

La arquitectura aprobada para el Capítulo \ref{chap:hibrido} prescinde del Proceso Gaussiano como predictor directo de biomasa —estadísticamente inviable con el volumen de datos destructivos disponible ($N \approx 30$) y metodológicamente ajena al objetivo específico 3 del proyecto, que exige una red neuronal— y adopta en su lugar un MLP como \emph{calibrador paramétrico} del DEB. El MLP opera sobre el subconjunto de primera sensibilidad $\{\tau_h, \rho_{\mathrm{conv}}, T_{\mathrm{opt}}\}$ emitiendo correcciones de escala $\boldsymbol{\alpha}_k$ respecto a los valores calibrados en este capítulo:
\begin{equation*}
    \theta_k^{\mathrm{final}} = \theta^{\mathrm{Cap3}} \odot (1 + \boldsymbol{\alpha}_k), \qquad \theta \in \{\tau_h, \rho_{\mathrm{conv}}, T_{\mathrm{opt}}\},
\end{equation*}
de modo que $\boldsymbol{\alpha}_k = \mathbf{0}$ recupera la calibración local del Capítulo \ref{chap:modelo}, mientras que $|\alpha_k| \ll 1$ representa ajustes inducidos por la variabilidad inter-replica. El MLP nunca ``ve'' biomasa durante el entrenamiento; su única supervisión es el residuo de \ch{CO2} $\varepsilon_k = s_k - J_{\ch{CO2}}^{\mathrm{DEB}}(\boldsymbol{\theta}_0)$, lo que evita el problema inverso mal condicionado y preserva la interpretabilidad fisiológica de cada parámetro.

\paragraph{Delimitación del alcance.} El metano (\ch{CH4}) fue excluido explícitamente de las variables de estado del modelo DEB, dado que su generación implica rutas fermentativas anaeróbicas no contempladas en el metabolismo aeróbico larval. En el Capítulo \ref{chap:metricas}, el \ch{CH4} se recupera como variable residual de diagnóstico: su detección por el sensor MQ-4 señala desviaciones del régimen aeróbico que el modelo DEB no predice, activando alertas operativas de anaerobiosis. Esta arquitectura —modelo mecanicista para la tendencia base, sensor de \ch{CH4} para las anomalías— preserva la parsimonia del DEB sin sacrificar la utilidad diagnóstica del sistema.

En síntesis, este capítulo ha transformado los datos crudos validados del Capítulo \ref{chap:instrumentacion} en una narrativa bioenergética cuantitativa interpretable: cada gramo de \ch{CO2} medido por el sensor NDIR puede ahora asociarse con un flujo metabólico específico (mantenimiento, síntesis de biomasa o síntesis de lípidos) bajo un conjunto de parámetros y condiciones ambientales explicitadas. El modelo DEB escalado a reactor batch, con sus ecuaciones, parámetros semilla, restricciones físicas y subconjunto de primera sensibilidad, queda listo para ser calibrado, validado e integrado en el sistema híbrido del Capítulo \ref{chap:hibrido}.
